% Contrato UC-14: un tema por fichero; contenido introductorio para validar la carga.
\bloque{funciones}
\tema{dominio}{Dominio de una función}
\dificultad{1}
\section{Dominio y valores de una función}
Una función relaciona cada elemento de su dominio con un único valor de salida.
Para estudiar una función conviene identificar primero qué entradas están permitidas.
El dominio depende de las operaciones que aparecen en la expresión y del contexto
del problema. En una fracción el denominador debe ser distinto de cero. En una
raíz cuadrada real el radicando debe ser mayor o igual que cero. Estas condiciones
se comprueban conjuntamente cuando una expresión contiene varias operaciones.
Por ejemplo, la expresión $f(x)=\frac{1}{x}$ tiene sentido real para valores
positivos y negativos, pero no está definida cuando el argumento vale cero.
La gráfica representa pares de entrada y salida, y ayuda a observar dónde cambia
el comportamiento. Una interrupción de la curva puede corresponder a una entrada
excluida del dominio; nunca debe unirse automáticamente a través de ese punto.
Antes de calcular valores, escribe las restricciones y comprueba cada una.
Después elige entradas sencillas a ambos lados de los puntos excluidos y compara
sus imágenes. Explica siempre por qué una entrada pertenece al dominio.
